\chapter{相关工作}
\label{chap:related_work}

本章系统回顾与本文最相关的研究脉络，从RLFT的基本范式演进、采样效率优化、难度感知预算分配、历史样本复用，到系统层优化等多个维度，梳理现有研究的进展与空白。

\section{大语言模型强化学习后训练与可验证奖励范式}
\label{sec:rlft_evolution}

\subsection{从监督微调到RLHF再到RLVR}
\label{subsec:sft_to_rlvr}

大语言模型的训练通常经历预训练、监督微调与后训练三个阶段。其中，预训练负责学习一般性的语言统计规律与知识表示，监督微调进一步塑造模型的指令跟随能力，而后训练则通过偏好信号、奖励信号或环境反馈，对模型行为进行面向目标场景的再优化。在复杂推理任务中，尤其是在数学、代码与形式逻辑等"结果可验证"场景中，强化学习后训练（Reinforcement Learning for Fine-Tuning/Post-Training, RLFT）已经成为提升模型推理能力的重要技术路线。

\textbf{监督微调的局限性。} 在这一演进链条中，监督微调（SFT）负责把预训练模型从通用语言建模任务迁移到指令跟随场景。然而，其根本机制仍然是对高质量示范分布的模仿，因此对数据质量与覆盖范围高度敏感。对于数学、代码与复杂推理任务而言，SFT面临两个突出问题：一是高质量过程监督数据昂贵且稀缺，二是即便拥有充足示范，模型学到的往往仍是"像训练集一样回答"，而不是稳定可迁移的推理策略。这一局限导致SFT方法在分布外问题上性能下降明显。

\textbf{RLHF的突破与瓶颈。} 为了突破纯监督模仿的局限，RLHF进一步将后训练形式化为基于人类偏好的策略优化问题。其经典流程通常先以SFT初始化策略模型，再用人类比较数据训练奖励模型，最后通过PPO等策略优化算法提升被偏好输出的概率。RLHF在通用对话、安全性与帮助性等维度取得了显著成功。然而，其上限又受制于奖励模型本身的质量，容易受到偏好噪声、数据规模限制与reward hacking的影响。对于具有客观评价标准的推理任务而言，人类偏好模型可能不是最优选择。

\textbf{RLVR的优势。} 针对这些问题，基于可验证奖励的强化学习（RLVR）提供了更适合推理任务的替代路径。RLVR不再依赖主观偏好模型，而是直接调用可执行验证器、答案比对器、单元测试或符号规则，对模型输出进行自动判分。这一路线的代表性成果（如DeepSeek-R1-Zero）已经证明，仅依赖结果可验证奖励，也能在数学推理等任务上显著激发模型潜在能力。

与传统监督学习不同，强化学习后训练不要求模型给出唯一标准的中间推理链条，而是允许模型围绕同一输入生成多个候选回答，并利用外部验证器对结果进行自动打分。对于数学任务，这一验证器通常由答案比对规则、程序执行器或符号计算器构成；对于代码任务，则可由单元测试、执行结果或形式约束进行评估。这一范式通常被称为可验证奖励强化学习（Reinforcement Learning with Verifiable Rewards, RLVR），其优势在于无需人工逐条构造细粒度监督信号，能够直接围绕最终任务目标进行优化。

\textbf{主流优化算法。} 在优化方法上，当前大语言模型后训练广泛采用基于策略梯度的近端更新框架，例如PPO、RLOO及其针对大语言模型场景的变体。其中，Group Relative Policy Optimization（GRPO）由于无需显式价值网络、实现简洁、与结果可验证任务适配性较强，而成为当前推理类后训练中的代表性方法之一。GRPO对每个Prompt采样固定数量的回答，并根据组内奖励差异构造相对优势。该思路在工程上具有较高可扩展性，但也把训练信号是否存在这一问题，直接绑定到了有限样本组内是否能观察到奖励差异之上。

\subsection{关于强化学习能力边界的争论}
\label{subsec:rl_capability_debate}

围绕RL是否真正"创造"了新的推理能力，近期出现了具有代表性的学术讨论。一些研究系统比较了多类RLVR算法与对应基座模型在大$k$条件下的Pass@k表现，并指出当前RLVR更多是在低采样预算下提升了正确推理路径被触发的概率，而非显著突破基座模型原有能力边界。这一结论虽然并不否定RLFT的价值，却对"采样是否足够充分"提出了更尖锐的问题：如果RL的主要作用是激发基座模型已具备但未充分显露的能力，那么采样效率就不再只是训练加速问题，而会直接决定这些潜在能力是否能被训练过程观察到、放大并稳定固化。

这一观点触发了后续研究的两类重要跟进。其一，一些工作承认现有post-training可能面临能力天花板，并尝试从训练状态建模层面重新打开能力扩张空间。其二，另一些工作试图调和"边界收缩"与"边界扩张"两种看法，提出更具阶段性的动态解释框架。这些研究的共同意义在于，它们都把注意力重新引向一个被长期低估的问题：\textbf{训练过程能否真正观察到模型已有但低概率出现的正确路径}。

\section{针对信号坍塌的采样效率优化}
\label{sec:signal_efficiency_optimization}

\subsection{早期应对策略}
\label{subsec:early_solutions}

面对信号坍塌问题，现有研究提出了若干直接的应对思路：

\textbf{被动过滤（Passive Filtering）} 在完成rollout之后，直接删除那些组内奖励完全一致、无法构造有效优势的Prompt，仅利用剩余样本更新模型。该方法虽然避免了零梯度样本直接进入更新，但会引入训练分布偏移。随着模型能力提升，简单题更容易因正向坍塌被剔除，难题则更容易因负向坍塌被剔除，最终保留下来的往往是中等难度样本。此外，被动过滤并不减少已发生的推理成本。

\textbf{固定大规模采样（Large Fixed-Budget Sampling）} 通过将采样数$n$增大至64、128乃至512来降低坍塌概率。该方法在理论上可以逼近充分采样的极限，但计算代价极高。在8张H100 GPU的训练配置下，将$n$从8提升至128会使单步训练时间由约3分钟增至接近50分钟。这对于实际项目而言是不可承受的。

\textbf{两阶段预算分配（Two-stage Budget Allocation）} 先使用少量样本估计每个Prompt的难度或通过率，再据此分配后续采样预算。然而这类方法面临三个困难：前置估计本身需要额外样本；小样本估计误差较大；离线或准离线的估计流程不易与大规模在线训练系统紧密耦合。

总体而言，这三类早期思路已经揭示了一个基本事实：固定rollout数并非最优设计，采样预算需要围绕"哪里更可能产生学习信号"进行重新组织。但它们仍然要么停留在采样之后的被动补救，要么依赖高成本的大样本估计。

\subsection{在线自适应采样}
\label{subsec:online_adaptive_sampling}

为了克服上述局限，近期研究转向\textbf{在线自适应采样}。这一方向的核心思想是：不再为每个Prompt固定分配相同数量的回答，而是在采样过程中根据已观测到的奖励结构，动态决定是否继续追加预算。

Reinforce-Ada是这一方向中极具代表性的工作。该方法通过successive elimination式的在线过程，将估计与采样交织进行。它在每一轮仅对仍处于活跃状态的Prompt继续采样，并在达到特定退出条件（如Balance条件：同时观测到至少一个正样本和一个负样本）后停止对该Prompt的进一步预算投入。其重要贡献在于明确指出：GRPO等固定组大小方法中的信号缺失并不一定意味着Prompt本身没有可学习信息，很多时候只是小样本下未能观测到正负奖励差异而已。

然而，现有在线自适应采样研究也存在若干未被充分解决的问题。首先，许多方法仍默认不同难度Prompt应遵循相同的退出逻辑，没有追问：在相同计算成本下，简单题与难题的单位信息收益是否真的相同。其次，多轮采样在统计上虽然更高效，但其常见实现方式往往忽略了现代推理引擎中的KV Cache复用特性。再次，当采样规模在不同轮次间动态变化时，不同轮次样本的探索强度是否需要同步调整，以及由此带来的分布偏差如何校正，也尚未形成系统解法。

\subsection{难度感知与数据中心化选择}
\label{subsec:difficulty_aware_selection}

在自适应采样方向进一步发展后，研究者开始意识到，仅仅做到"有信号就继续、没信号就停止"仍然不够，因为不同Prompt的信息密度并不相同。由此，一批工作转向更细粒度地研究\textbf{难度感知}、\textbf{能力对齐}与\textbf{数据中心化预算分配}问题。

\textbf{基于难度的动态rollout分配。} 一些工作指出现有RL方法为所有问题分配相同回答数并不高效。简单问题继续投入额外rollout通常收益有限，而困难问题则需要更多样本才能暴露出潜在正确路径。相应的方法依据题目难度动态调整rollout budget，并同时引入动态温度调节与退火机制，以维持训练过程中的策略熵。

\textbf{能力-难度对齐。} 进一步的工作指出，真正决定样本训练价值的，不仅是问题难度本身，还包括其与当前模型能力之间的匹配关系。样本价值由"样本难度与模型当前可学习区间之间的相对位置"共同决定，这与传统curriculum learning在理念上相通。

\textbf{从梯度方差到信息效率。} 另一类工作从梯度方差最小化角度给出rollout分配原则，在固定计算预算下求解最优rollout分配。这为后续关于单位成本信息效率的讨论提供了更直接的统计学支点。

\section{历史样本复用与系统级优化}
\label{sec:sample_reuse_and_system}

\subsection{响应复用与前缀复用}
\label{subsec:response_prefix_reuse}

除了决定"采什么"和"采多少"，近期研究还关注：\textbf{已经生成过的样本是否必须一次性使用后丢弃}。在标准在线RLFT中，rollout往往被默认视为临时样本，即本轮用于构造优势与更新策略后便不再复用。然而，大语言模型rollout的生成成本极高，若历史样本中仍保留有价值信号，则简单丢弃显然会造成额外浪费。

一类工作通过\textbf{响应复用（Response Reuse）}将历史上已经生成过的正确回答重新纳入训练信号构造过程。与仅依赖当前步样本的策略相比，这有效缓解了GRPO中组内奖励完全一致导致的梯度消失问题，并显著降低了反复执行rollout的成本。对于结果可验证任务而言，"曾经采到过的正样本"本身就是高价值资产，不应被当前训练步之外的时间边界所束缚。

更细粒度的\textbf{前缀复用（Prefix-level Reuse）}方向进一步指出，同一查询下的多条rollout往往在早期推理步骤上高度相似。若每次都从头独立生成，会造成显著的前缀重复计算。前缀复用通过缓存并共享这些早期推理段，从而在数据效率与系统效率之间建立直接联系。这与本文关注的KV Cache复用问题具有紧密的相关性。

\subsection{系统效率视角}
\label{subsec:system_efficiency}

从系统执行角度看，大语言模型在线采样通常包含Prefill与Decode两个阶段：前者负责处理输入Prompt并构建KV Cache，后者逐token生成回答并持续读写缓存状态。对于同一Prompt下的多条候选回答而言，其前缀通常完全相同，因此从理论上讲只需执行一次Prefill，随后由多条回答共享同一前缀缓存即可。这一前缀共享机制已经成为现代推理引擎（vLLM、SGLang等）的重要优化基础。

\textbf{并发控制与部分rollout。} 在同步式RLFT系统中，极长轨迹会拖慢整批样本返回，使得训练不得不等待少数长尾回答完成，造成GPU空转与吞吐下降。一类工作（如CoPRIS）引入concurrency-controlled partial rollout：系统维持固定数量的并发rollout，当收集到足够样本后即可提前终止本轮，并把未完成轨迹缓冲到后续阶段继续使用。由于这些轨迹在时间上跨越了策略更新边界，相应的方法进一步通过importance sampling correction对stale trajectory带来的分布偏差进行修正。

\textbf{生成时在线决策。} 另一类工作甚至把优化时机推进到了生成过程中。在解码尚未结束时就使用轻量级质量头预测partial rollout的成功概率，并对低价值候选进行早停剪枝。与训练后筛除低价值样本相比，这能够直接减少无效token的生成，并通过inference engine内部重排批次，提高后续计算的整体效率。

这一系列工作表明：系统效率不应被视为与算法无关的工程细节，而应成为采样框架设计的内生约束。即使某种采样策略在统计上更优，若其执行方式无法充分利用底层系统特性，也可能无法转化为真实训练收益。

\section{现有研究的总体比较与空白}
\label{sec:research_gaps}

综合上述文献，可以把现有研究大致归纳为四条主线：

\begin{enumerate}
    \item \textbf{固定组优化下的被动补救}：包括过滤无信号样本与简单扩大rollout数。这类方法实现直接，但大多停留在"信号缺失之后如何减轻损失"。
    
    \item \textbf{在线自适应采样}：代表性工作如Reinforce-Ada，其核心是根据当前已观测信号动态决定是否继续采样。该方向首次系统性地把"信号能否被采到"当作主要研究对象。
    
    \item \textbf{难度感知与数据中心化选择}：这类工作普遍承认样本并不等价，尝试依据难度、能力匹配程度、策略置信状态或样本信息价值进行更精细的预算调度。
    
    \item \textbf{历史样本复用与系统级优化}：分别从响应复用、前缀复用、部分rollout缓存和生成时剪枝等角度，提升历史样本利用率与执行吞吐。
\end{enumerate}

尽管现有研究已经极大推进了RLFT采样效率问题的认识，但在本文所关注的场景下，仍至少存在以下\textbf{三个尚未被充分解决的空白}：

\begin{enumerate}
    \item \textbf{理论与实证的不对称}：现有在线自适应采样通常默认简单题与难题遵循相同的退出逻辑，却较少直接从"单位成本信息效率"角度比较不同难度区域的真实预算价值。
    
    \item \textbf{算法层与系统层的耦合不足}：现有采样工作主要关注统计效率，较少把KV Cache复用、批次组织方式与推理引擎并行模式纳入方法设计。多轮固定批次追加的实现方式与现代大语言模型推理系统的前缀共享机制之间仍存在显著错配。
    
    \item \textbf{动态探索与分布校正的统一不足}：当采样规模在不同轮次间动态变化时，探索强度与分布偏差的联动问题尚未形成完整解决方案。一部分工作认识到温度调节对探索的重要性，另一部分工作强调importance sampling对跨分布样本校正的必要性，但把两者统一为一个框架的工作仍然较少。
\end{enumerate}

从这一意义上说，现有研究虽然已经证明"固定组大小不是最优解"，但对于\textbf{信息效率、系统效率与统计一致性三者如何协同设计}仍缺乏统一视角。本文正是试图在这一交汇点上继续推进。

\section{本章小结}
\label{sec:related_work_summary}

本章回顾了与本文最相关的研究脉络。总体来看，围绕大语言模型强化学习后训练中的采样效率问题，已有研究已从早期的被动过滤与暴力大采样，逐步发展到在线自适应采样、难度感知预算分配、能力-难度对齐、历史样本复用以及系统级执行优化等多个方向。

这些工作共同说明：在RLFT中，真正稀缺的并不是"样本数量"本身，而是在有限预算下能够形成有效策略更新的\textbf{高价值信号}。然而，现有工作仍未充分回答：不同难度样本在单位成本下的信息产出是否等价；采样批次组织方式能否与推理系统KV Cache复用机制协同优化；以及当采样规模动态变化时，探索强度与策略分布校正应如何联合设计。

围绕这些问题，下一章将给出本文的完整方法框架。
